\section*{Chapitre \ref{ch:m1:zero-day-options-and-retail-flow} --- Options à zéro jour, options hebdomadaires et flux d'options de particuliers}

\begin{solution}{exo:m1:zero-day-options-and-retail-flow:1}
$0.8 \times 16\%/\sqrt{252} = 0.81\%$ de l'indice~; à 32~\%, 1.61~\%. (La
valeur exacte à 16~\% est de 0.80~\%.)
\end{solution}

\begin{solution}{exo:m1:zero-day-options-and-retail-flow:2}
Couverture~: vente à découvert de $0.50 \times 500 \times 100 = 25\,000$ actions. Estimation linéaire pour
$+1\%$ (un dollar)~: \omterm{def:m1:zero-day-options-and-retail-flow:gamma}{delta} augmente de 0.396, soit $0.396 \times 50\,000 =
19\,800$ actions (19\,790 avant arrondi). Le desk est long \omterm{def:m1:zero-day-options-and-retail-flow:gamma}{gamma}~: il \emph{vend}
les actions à mesure que leur cours monte.
\end{solution}

\begin{solution}{exo:m1:zero-day-options-and-retail-flow:3}
$\sqrt{21} = 4.6$~; $\sqrt{390/15} = 5.1$. Le \omterm{def:m1:zero-day-options-and-retail-flow:gamma}{gamma} est multiplié par environ vingt-trois
entre un mois et le dernier quart d'heure.
\end{solution}

\begin{solution}{exo:m1:zero-day-options-and-retail-flow:4}
Le \omterm{def:m1:zero-day-options-and-retail-flow:gamma}{delta} passe de 0.50 à 0.84~: vendre $0.34 \times 50\,000 = 16\,900$ actions
(16\,873 exactement), contre 19\,790 selon l'estimation linéaire. Le \omterm{def:m1:zero-day-options-and-retail-flow:gamma}{gamma} est la pente
du \omterm{def:m1:zero-day-options-and-retail-flow:gamma}{delta} à 100~; à 101, un dollar représentant 1.0 écart-type de la
journée restante, l'option est déjà bien dans la monnaie et son \omterm{def:m1:zero-day-options-and-retail-flow:gamma}{gamma} a
diminué~: le \omterm{def:m1:zero-day-options-and-retail-flow:gamma}{delta} sature à un.
\end{solution}

\begin{solution}{exo:m1:zero-day-options-and-retail-flow:5}
$2.3 \times 10^6 \times 100 \times 6\,000 = \$1.38$ billion de notionnel. Les flux
de couverture dépendent du \omterm{def:m1:zero-day-options-and-retail-flow:gamma}{delta} et du \omterm{def:m1:zero-day-options-and-retail-flow:gamma}{gamma}, non du notionnel~: la plupart de ces options sont hors de la
monnaie, avec de petits \omterm{def:m1:zero-day-options-and-retail-flow:gamma}{deltas}, et un hedger ne négocie que la \emph{variation} du \omterm{def:m1:zero-day-options-and-retail-flow:gamma}{delta}.
Et c'est la position \emph{nette} des dealers qui est couverte~: les clients achètent et vendent
les mêmes strikes, les contrats sont ouverts et clôturés plusieurs fois au cours de la
journée, de sorte que la position nette ne représente qu'une faible fraction du volume. Deux ou trois ordres de
grandeur séparent les deux nombres.
\end{solution}

\begin{solution}{exo:m1:zero-day-options-and-retail-flow:6}
Short \omterm{def:m1:zero-day-options-and-retail-flow:gamma}{gamma}~: il achète des actions à mesure que le cours monte de 3~\% dans la dernière heure (le
\omterm{def:m1:zero-day-options-and-retail-flow:gamma}{delta} de ses calls vendus monte) et les vend lorsque le cours retombe le
lendemain matin. Il a acheté haut et vendu bas~: une perte de couverture à peu près proportionnelle au
\omterm{def:m1:zero-day-options-and-retail-flow:gamma}{gamma} multiplié par le carré du mouvement. En contrepartie, il a encaissé la prime des options, c'est-à-dire la
\omterm{def:m1:volatility-first-contact:iv}{volatilité implicite}~; il gagne si les mouvements réalisés qu'il doit
suivre sont plus faibles que ceux que la prime rémunère.
\end{solution}

\begin{solution}{exo:m1:zero-day-options-and-retail-flow:7}
0.77, 1.00 et 1.42. Pour $r_t = \varepsilon_t + f\,r_{t-1}$, la somme des rendements
sur un long horizon vaut environ $\sum\varepsilon_t/(1-f)$, donc le rapport vaut
$1/(1-f)$~: 0.77 pour $f = -0.3$ et 1.43 pour $f = +0.3$. L'asymétrie est
celle de $1/(1-f)$~: le rapport diverge quand $f \to 1$.
\end{solution}

\begin{solution}{exo:m1:zero-day-options-and-retail-flow:8}
(i) Les options du jour même, ouvertes et clôturées pendant la séance, ne figurent pas du tout dans les
\omterm{def:m1:futures-contracts-and-exchanges:oi}{positions ouvertes} de la veille~: le chiffre décrit les positions à plus longue échéance.
(ii) Le signe est supposé~: si les clients vendent des calls (vente couverte) et vendent aussi
des puts (encaissement de la prime), les dealers sont longs des deux et le chiffre réel est
plus élevé~; si les clients achètent des calls, comme le font les particuliers, les dealers sont vendeurs de ces calls et
le signe s'inverse. (iii) Le \omterm{def:m1:zero-day-options-and-retail-flow:gamma}{gamma} à l'approche de l'échéance dépend fortement de la position de l'indice
par rapport à chaque strike et de l'heure~: un chiffre calculé à 09:00 pour
«~aujourd'hui~» est périmé dès 10:00. On peut ajouter que les flux de clients sur ces options
sont souvent dans les deux sens, de sorte que la position nette des dealers est bien plus petite que
ne le laissent supposer les \omterm{def:m1:futures-contracts-and-exchanges:oi}{positions ouvertes} ou le volume.
\end{solution}

\begin{solution}{pb:m1:zero-day-options-and-retail-flow:1}
\textbf{1.} 2 millions d'actions, 1.2~milliards de dollars.
\textbf{2.} Position short de 20\,000 calls de \omterm{def:m1:zero-day-options-and-retail-flow:gamma}{delta} 0.50~: position longue de 1 million d'actions.
\textbf{3.} $16\% \times \sqrt{60/98\,280} = 0.40\%$.
\textbf{4.} $0.168 \times 6 \times 2\,000\,000 = 2.02$ millions d'actions pour 1~\%.
Short \omterm{def:m1:zero-day-options-and-retail-flow:gamma}{gamma}~: le dealer \emph{achète} quand le cours monte et vend quand il baisse.
\textbf{5.} L'estimation dépasse la totalité du notionnel, ce qui est impossible~: le \omterm{def:m1:zero-day-options-and-retail-flow:gamma}{delta} ne peut
passer que de 0.50 à 1. Le chiffre linéaire n'est valable que pour des mouvements petits devant
un écart-type, 0.40~\%, et un mouvement de 1~\% en représente deux et demi.
\textbf{6.} $0.50 + 472\,000/2\,000\,000 = 0.74$ et $0.50 + 987\,000/2\,000\,000 =
0.99$.
\textbf{7.} Après $+0.25\%$, soit six dixièmes d'écart-type, le \omterm{def:m1:zero-day-options-and-retail-flow:gamma}{delta} a
déjà parcouru la moitié de sa plage restante~; au-delà d'un écart-type
il reste peu à couvrir.
\textbf{8.} $2.02/0.792 = 2.55 = \sqrt{390/60}$; $4.04/2.02 = 2.0 =
\sqrt{60/15}$.
\textbf{9.} Une heure moyenne échange $60/6.5 = 9.2$ millions d'actions~: le chiffre
linéaire représente 22~\% de ce total, pour un seul dealer et un seul strike.
\textbf{10.} À 600.30, les calls sont tout juste dans la monnaie et le dealer est long de
l'essentiel des 2 millions d'actions~; à 599.70, il devrait en détenir peu. À cinq
minutes de la fin, chaque franchissement du strike le conduit à vendre, puis à racheter, une
grande partie des 2 millions d'actions.
\textbf{11.} Il vend près de 599.70 et achète près de 600.30~: jusqu'à 60 cents sur les
actions qu'il échange, par oscillation, auxquels s'ajoutent la fourchette et l'impact. C'est la
\omterm{def:m1:volatility-first-contact:realised}{volatilité réalisée} dont il est short, payée en espèces.
\textbf{12.} La prime des 20\,000 calls, versée par les clients qui
les avaient achetés ce matin~: environ 0.4~\% de 1.2~milliards de dollars, soit environ 5~millions de dollars, si elles sont vendues à
l'ouverture.
\textbf{13.} Assignation~: si le fonds clôture à 600 ou à proximité, il ne sait pas
combien de calls seront exercés, donc combien de ses actions de couverture il doit
livrer (voir le \cref{pb:m1:option-contracts-and-exchanges:1}).
\textbf{14.} Il vendrait au-dessus de 600 et achèterait en dessous~: sa couverture fournit
de la liquidité des deux côtés du strike et tend à retenir le cours à ce niveau, le
pinning du \cref{pb:m1:option-contracts-and-exchanges:1}, question 16.
\textbf{15.} Elles pourraient l'être si les dealers étaient fortement short \omterm{def:m1:zero-day-options-and-retail-flow:gamma}{gamma} près de la
monnaie~: leurs achats sur les hausses et leurs ventes sur les baisses, croissant à mesure que
l'échéance approche, s'auto-alimenteraient (voir la \cref{fig:m1:zero-day-options-and-retail-flow:feedback}).
Ce qui le limite, c'est que les clients achètent et vendent à la fois ces options, de sorte que
le \omterm{def:m1:zero-day-options-and-retail-flow:gamma}{gamma} net des dealers ne représente qu'une faible fraction du volume, et de l'un ou l'autre signe.
\textbf{16.} Elle est mesurée une fois par jour à partir des positions de la veille~;
les opérations du jour même s'ouvrent et se clôturent dans la séance.
\textbf{17.} De nombreux petits gains et de rares pertes, plusieurs fois la prime~: une
distribution à longue queue gauche. La taille des positions doit découler de la perte
subie un jour où l'indice bouge de trois ou quatre écarts-types journaliers, et non de
la moyenne.
\textbf{18.} Volume~: chaque échéance est un nouveau produit qui n'exige aucune nouvelle
cotation, attire la demande des fonds event-driven et des particuliers, et est réservée à la
bourse qui détient la licence sur l'indice.
\textbf{19.} \textbf{2.02 millions d'actions pour 1~\%, soit 1.2~milliards de dollars
en actions, soit la totalité du notionnel.}
\textbf{20.} Son \omterm{def:m1:zero-day-options-and-retail-flow:gamma}{delta} n'est plus un nombre qui évolue lentement, mais un interrupteur que
le sous-jacent peut faire basculer.
\end{solution}

\begin{solution}{iq:m1:zero-day-options-and-retail-flow:1}
Le \omterm{def:m1:zero-day-options-and-retail-flow:gamma}{gamma} est la vitesse de variation du \omterm{def:m1:zero-day-options-and-retail-flow:gamma}{delta} par rapport au sous-jacent. \omterm{def:m1:zero-day-options-and-retail-flow:gamma}{Gamma} long
(options achetées)~: le \omterm{def:m1:zero-day-options-and-retail-flow:gamma}{delta} croît avec le cours, donc le hedger vend à la hausse et
achète à la baisse, gagnant sur le mouvement et payant la décroissance temporelle. \omterm{def:m1:zero-day-options-and-retail-flow:gamma}{Gamma} court~:
l'inverse~; le hedger court après le marché et est rémunéré par la prime pour cela.
\iqlookfor{le sens des transactions de couverture et le compromis
entre décroissance temporelle et mouvement.}
\end{solution}

\begin{solution}{iq:m1:zero-day-options-and-retail-flow:2}
\omterm{def:m1:volatility-first-contact:atm}{À la monnaie}, il croît comme $1/\sqrt T$ et devient infini à l'échéance~: le \omterm{def:m1:zero-day-options-and-retail-flow:gamma}{delta}
passe brusquement de zéro à un au niveau du prix d'exercice. Hors de la monnaie, il croît d'abord,
à mesure que la distribution se resserre vers le voisinage du prix d'exercice, puis
s'effondre vers zéro dès que le prix d'exercice est à plusieurs $\sigma\sqrt T$ de distance. Le \omterm{def:m1:zero-day-options-and-retail-flow:gamma}{gamma}
se concentre dans une bande de largeur $\sigma\sqrt T$ autour du prix d'exercice.
\iqlookfor{les deux comportements et la bande.}
\end{solution}

\begin{solution}{iq:m1:zero-day-options-and-retail-flow:3}
$0.8\,\sigma\sqrt T$~: avec une volatilité de 16~\% par an, un jour vaut $16/\sqrt{252} \approx 1\%$,
donc le straddle coûte environ 0.8~\% de l'indice, soit 48 points à 6\,000. Il atteint le point mort
si l'indice clôture à plus de 0.8~\% de distance, soit environ 0.8
écarts-types journaliers, ce qui correspond à une probabilité d'environ 42~\%.
\iqlookfor{la règle du seize et une probabilité.}
\end{solution}

\begin{solution}{iq:m1:zero-day-options-and-retail-flow:4}
Classer chaque transaction d'options comme achat ou vente du client à partir du ruban (codes de capacité
lorsqu'ils existent, sinon règle de cotation et taille), cumuler les positions côté dealers par prix d'exercice
au fil de la journée, et réévaluer le \omterm{def:m1:zero-day-options-and-retail-flow:gamma}{gamma} au spot et au temps courants. Test~:
régresser la \omterm{def:m1:volatility-first-contact:realised}{volatilité réalisée} ultérieure, ou l'autocorrélation des rendements à court horizon,
sur cette estimation, hors échantillon, en contrôlant la \omterm{def:m1:volatility-first-contact:iv}{volatilité implicite} et l'heure de la journée~;
vérifier que le signe s'inverse lorsque l'estimation s'inverse. Être
prêt à constater que l'effet est faible.
\iqlookfor{transactions signées plutôt que \omterm{def:m1:futures-contracts-and-exchanges:oi}{positions ouvertes}~; un test réfutable.}
\end{solution}

\begin{solution}{iq:m1:zero-day-options-and-retail-flow:5}
Racheter les options et payer le solde de la prime~: la sortie propre. Acheter
d'autres options aux prix d'exercice voisins pour plafonner le \omterm{def:m1:zero-day-options-and-retail-flow:gamma}{gamma}. Continuer à couvrir, en acceptant le
coût de chaque franchissement, avec une bande de couverture plus large pour éviter les allers-retours.
Ou cesser de couvrir et porter l'issue binaire, dimensionnée pour rester supportable. Le pire choix
est de couvrir en continu avec une bande étroite.
\iqlookfor{la clôture par défaut~; les bandes~; le dimensionnement.}
\end{solution}

\begin{solution}{iq:m1:zero-day-options-and-retail-flow:6}
Mécanisme~: les dealers nets vendeurs de \omterm{def:m1:zero-day-options-and-retail-flow:gamma}{gamma} du jour amplifient les mouvements, les dealers nets
acheteurs les amortissent~; le \omterm{def:m1:zero-day-options-and-retail-flow:gamma}{gamma} par contrat est énorme près du prix d'exercice. Contre un effet
de grande ampleur~: les flux clients sont bilatéraux et une grande partie consiste en spreads, si bien que le
\omterm{def:m1:zero-day-options-and-retail-flow:gamma}{gamma} net des dealers est faible~; les positions sont soldées à la clôture~; la couverture se fait
sur des contrats à terme très liquides. Données nécessaires~: le \omterm{def:m1:zero-day-options-and-retail-flow:gamma}{gamma} net des dealers par minute, à partir des
transactions signées~; son lien avec l'autocorrélation des rendements intrajournaliers et la volatilité~;
des études d'événements les jours de flux unilatéraux importants~; la comparaison de la volatilité
intrajournalière avant et après l'introduction des échéances quotidiennes, en contrôlant le
niveau de la volatilité.
\iqlookfor{net contre brut~; les bonnes données.}
\end{solution}
